\section{Механизм управления дивергентным потоком}

\subsection{Теория PDOM-реконвергенции}

\subsubsection{Граф управляющего потока и постдоминаторы}

Граф управляющего потока (CFG) программы определяется как:
\begin{equation}
  G = (V,\, E,\, \mathit{entry},\, \mathit{exit}),
\end{equation}
где $V$~--- множество базовых блоков,
$E \subseteq V \times V$~--- рёбра передачи управления,
$\mathit{entry}$ и $\mathit{exit}$~--- специальные узлы.

\textbf{Определение (постдоминатор).} Узел $D$ является постдоминатором
узла $N$ ($D \;\mathrm{pdom}\; N$), если каждый путь из $N$ в
$\mathit{exit}$ проходит через $D$.

\textbf{Определение (непосредственный постдоминатор).}
Непосредственный постдоминатор узла $N$:
\begin{equation}
  \mathrm{IPDOM}(N) = \operatorname*{argmin}_{D:\; D \;\mathrm{pdom}\; N,\; D \neq N}
    |E^*[N, D]|,
\end{equation}
то есть ближайший постдоминатор (за минимальное число шагов).

\textbf{Теорема.} $\mathrm{IPDOM}(N)$ существует для каждого узла
$N \neq \mathit{exit}$.

\textit{Доказательство} строится по индукции по структуре CFG:
сначала для ациклических CFG, затем обобщается на произвольные CFG
методом фиксированной точки~\cite{Cytron1991}.

\subsubsection{Теорема о корректности PDOM-реконвергенции}

\textbf{Теорема (PDOM-корректность).}
Пусть при исполнении условного перехода $N$ варьируется
активность потоков: $M_{\mathrm{taken}} \neq 0$ и
$M_{\mathrm{not}} \neq 0$. Если точкой реконвергенции выбирается
$\mathrm{IPDOM}(N)$, то:
\begin{enumerate}
  \item Все потоки варпа воссоединятся в $\mathrm{IPDOM}(N)$.
  \item Ни один поток не пропустит инструкцию своего активного пути.
  \item После реконвергенции маска $M$ восстанавливается в полное значение.
\end{enumerate}

\textit{Обоснование.}

\textbf{Утверждение~1} (воссоединение).
Введём обозначения: \textit{ветка перехода}~--- подмножество потоков
с $P=1$, которые переходят на метку $\$L$; \textit{ветка продолжения}~---
подмножество потоков с $P=0$, которые продолжают выполнение с $\mathrm{PC}+1$.
По определению постдоминатора каждый путь из $N$ в $\mathit{exit}$
проходит через $\mathrm{IPDOM}(N)$.
Ветка перехода ($N \to \$L \to \cdots$) и ветка продолжения
($N \to \mathrm{PC}+1 \to \cdots$)~--- оба пути ведут в
$\mathit{exit}$, следовательно, оба обязательно пройдут через
$\mathrm{IPDOM}(N)$. Значит, все потоки встретятся в этой точке.

\textbf{Утверждение~2} (полнота путей). Потоки с маской $M_{\mathrm{taken}}$
исполняются только внутри ветки перехода, а потоки с маской
$M_{\mathrm{not}}$~--- только внутри ветки продолжения.
Ни один поток не получает нулевую маску пока его ветка активна,
поэтому ни одна инструкция активного пути не пропускается.

\textbf{Утверждение~3} (восстановление маски). В стек дивергенции
помещается фрейм реконвергенции с исходной маской $M$ и
$\mathrm{PC} = \mathrm{IPDOM}(N)$. Когда счётчик команд достигает этой
точки, алгоритм CheckStack извлекает фрейм и восстанавливает $M$.

Поскольку выбирается \textit{непосредственный} постдоминатор~---
ближайший к $N$~--- последовательное исполнение двух ветвей
минимально: потоки воссоединяются как можно раньше.

\subsection{Структура данных стека дивергенции}

\subsubsection{Тип \texttt{StackFrame} и семантика полей}

Каждый фрейм стека дивергенции описывается структурой:

\begin{lstlisting}[language=C++, caption={Структура StackFrame}]
struct StackFrame {
    uint64_t pc;
    uint32_t mask;
    bool     is_convergence;
};
\end{lstlisting}

Инварианты:
\begin{itemize}
  \item Если \texttt{is\_convergence~= true}: при достижении
        $\mathrm{PC} = \mathit{frame.pc}$ маска восстанавливается
        в \texttt{frame.mask} и фрейм снимается.
  \item Если \texttt{is\_convergence~= false}: дивергентный путь
        в очереди; при достижении его PC исполнение возобновляется
        с маской \texttt{frame.mask}.
\end{itemize}

Стек $S = [f_n, \ldots, f_1, f_0]$ (вершина~--- $f_n$):
фрейм реконвергенции всегда располагается \textit{ниже}
соответствующих дивергентных фреймов.

\subsubsection{Сложность операций и ограничение глубины}

Операции push/pop/top~--- O(1). При вложенности дивергенций
глубиной $k$: $|S| \leq 2k$.

Для типичных CUDA-ядер $k \leq 5$~--- $|S| \leq 10$.
Объём памяти: $\texttt{sizeof}(\texttt{StackFrame}) \cdot 10 = 160$~байт.

\subsection{Алгоритм обработки условного перехода}

\subsubsection{Алгоритм PDOM\_Diverge}

\begin{algorithm}[H]
  \caption{PDOM\_Diverge($M$, $P$, $\$L$, $\mathit{pc\_ipdom}$)}
  \label{alg:pdom_diverge}
  \KwIn{маска исполнения $M$, предикат $P$, метка перехода $\$L$,
        PC точки реконвергенции $\mathit{pc\_ipdom}$}
  \KwOut{обновлённые $M$, $\mathrm{PC}$, стек $S$}

  $M_{\mathrm{taken}} \leftarrow M \wedge P$\;
  $M_{\mathrm{not}} \leftarrow M \wedge \neg P$\;
  \textbf{push}$(S,\;\{\mathit{pc\_ipdom},\; M,\; \mathit{is\_conv} = \texttt{true}\})$\;
  \textbf{push}$(S,\;\{\mathrm{PC}+1,\; M_{\mathrm{not}},\; \mathit{is\_conv} = \texttt{false}\})$\;
  $M \leftarrow M_{\mathrm{taken}}$\;
  $\mathrm{PC} \leftarrow \mathit{offset}(\$L)$\;
\end{algorithm}

\textbf{Пояснение к шагам:}
\begin{itemize}
  \item Шаг~3 фиксирует точку воссоединения $\mathit{pc\_ipdom}$
        с исходной полной маской $M$.
  \item Шаг~4 откладывает ветку продолжения: при PC~$=\mathrm{PC}+1$
        она будет исполняться с маской $M_{\mathrm{not}}$.
  \item Шаги~5--6 немедленно переходят к исполнению ветки перехода.
\end{itemize}

\subsubsection{Цикл проверки вершины стека}

\begin{algorithm}[H]
  \caption{CheckStack($S$, $\mathrm{PC}$, $M$)}
  \label{alg:check_stack}
  \While{\textbf{not} $S.\mathit{empty}()$ \textbf{and} $S.\mathit{top}().pc = \mathrm{PC}$}{
    \eIf{$S.\mathit{top}().\mathit{is\_conv}$}{
      $M \leftarrow S.\mathit{top}().\mathit{mask}$\;
      $S.\mathit{pop}()$\;
    }{
      $M \leftarrow S.\mathit{top}().\mathit{mask}$\;
      $S.\mathit{pop}()$\;
      \textbf{break}\tcp*{начало дивергентного пути}
    }
  }
\end{algorithm}

Алгоритм CheckStack вызывается \textit{до} исполнения каждой инструкции.
Сложность: O($k$) за одну проверку $\Rightarrow$ общая сложность
O($k \cdot N_{\mathrm{instrs}}$) за исполнение ядра с $k$ уровнями дивергенции.

\subsection{Угловые случаи}

\subsubsection{Вложенные дивергенции: корректность при $k \geq 2$}

Рассмотрим пример: внешний дивергентный \texttt{if} и внутренний
дивергентный \texttt{if}. Пусть начальная маска $M_0 = \texttt{0xFFFFFFFF}$:

\begin{enumerate}
  \item Внешний \texttt{if}: PDOM\_Diverge создаёт $f_2$
        (реконвергенция, $M_0$) и $f_1$ (ветка «не взята», $M_{\mathrm{not\_outer}}$).
        Стек: $[f_2, f_1]$.
  \item Внутренний \texttt{if}: PDOM\_Diverge создаёт $f_4$
        (внутренняя реконвергенция) и $f_3$ (внутренняя ветка).
        Стек: $[f_2, f_1, f_4, f_3]$.
  \item Достижение внутренней реконвергенции: CheckStack снимает
        $f_4$ и $f_3$, стек: $[f_2, f_1]$.
  \item Достижение внешней реконвергенции: CheckStack снимает
        $f_2$ и $f_1$, стек: $[\,]$.
\end{enumerate}

LIFO-порядок снятия фреймов гарантирует восстановление масок
в порядке «внешнее-первым», то есть корректное поведение при любой
глубине вложенности.

\subsubsection{Дивергентные циклы}

В \texttt{while}-цикле с дивергентным условием выхода: часть потоков
достигает \texttt{bra~\$exit} раньше остальных. При каждой итерации
PDOM\_Diverge корректно разделяет «продолжающие» ($M_{\mathrm{continue}}$)
и «вышедшие» ($M_{\mathrm{exit}}$) потоки. Потоки, завершившие цикл,
не участвуют в последующих итерациях, так как их биты сброшены
в $M_{\mathrm{continue}}$.

\subsubsection{Ранний выход и переход к \texttt{isActive() = false}}

Инструкция \texttt{exit} при $M \neq 0$ обнуляет биты завершивших
потоков в маске:
$M' = M \setminus \{t \mid \text{поток } t \text{ исполняет exit}\}$.

При $M' = 0$ метод \texttt{isActive()} возвращает \texttt{false}:
диспетчер (Часть~2) останавливает варп. Часть~1 только обновляет
маску~--- сигнал останова передаётся через публичный метод.

\newpage
